\documentclass[11pt,a4paper]{article}
\usepackage[margin=25mm]{geometry}
\usepackage{amsmath,amssymb,amsthm,mathtools}
\usepackage[T1]{fontenc}
\usepackage{lmodern,microtype}
\usepackage[hidelinks]{hyperref}
\hypersetup{pdftitle={An explicit rank-two entropy-minimizing Fano degeneration},pdfauthor={Alec Kriebel}}
\newtheorem{theorem}{Theorem}
\newtheorem{lemma}[theorem]{Lemma}
\newtheorem{proposition}[theorem]{Proposition}
\newcommand{\C}{\mathbb C}
\newcommand{\Q}{\mathbb Q}
\newcommand{\Z}{\mathbb Z}
\newcommand{\R}{\mathbb R}
\newcommand{\ord}{\operatorname{ord}}
\newcommand{\Bl}{\operatorname{Bl}}
\newcommand{\gr}{\operatorname{gr}}
\newcommand{\Proj}{\operatorname{Proj}}
\newcommand{\MA}{\operatorname{MA}}
\newcommand{\Ric}{\operatorname{Ric}}
\newcommand{\Res}{\operatorname{Res}}
\title{An explicit rank-two entropy-minimizing\\ Fano degeneration}
\author{Alec Kriebel\\{\small Independent researcher}\\
{\small\href{https://orcid.org/0009-0001-9320-500X}{ORCID: 0009-0001-9320-500X}}}
\date{16 October 2026 \quad Version 1.0.0}
\begin{document}
\maketitle
\begin{abstract}
We construct a smooth Fano fivefold with a nondivisorial global
entropy-minimizing valuation and a nontrivial canonical first degeneration.
The example is the product of an explicit quadric blowup with the first
Hirzebruch surface. Its two positive valuation weights are specified by exact
moment equations; an elementary resultant argument, using the classical
Hermite--Lindemann theorem, proves that their ratio is irrational.
The associated graded anticanonical ring identifies the first central
variety, which is singular although the original fivefold is smooth.
Global optimality follows from the established optimal-degeneration and
soliton theorems after checking the product filtration and normalization.
The threefold factor and its optimal degeneration are due to Miao--Wang.
No claim of novelty or historical firstness is made.
\end{abstract}

\section{The construction and its target}

Blum's discussion of canonical Fano degenerations asks for explicit examples
whose first central variety differs from the original variety, and in
particular for one with a nondivisorial minimizing valuation
\cite[pp.~888--889]{OWR}. The construction below satisfies both requirements.
We work over \(\C\). By a divisorial real valuation we mean
\(c\,\ord_F\), where \(F\) is a prime divisor over the variety and \(c>0\)
is real. An irrational rescaling alone therefore does not prove
nondivisoriality.

Earlier explicit nontrivial first centers include Li--Li's
\(\mathrm{SO}_4\)-compactifications \cite[\S6]{LL}.
Their second example has a two-component optimizer; we have not verified
irrationality of its direction. Miao--Wang's divisorial threefold factor
already appears in their January 2024 first preprint.
We cite the revised edition for precise theorem numbering.
The present claim is an explicit construction, not a priority certificate.

In \(\mathbb P^4_{[u:x_0:x_1:x_2:x_3]}\), set
\[
 Q=\{u^2+x_0^2+x_1^2+x_2^2+x_3^2=0\},\qquad H=Q\cap\{u=0\},
\]
\[
 C=H\cap\{x_0^2+2x_1^2+3x_2^2+4x_3^2=0\},\qquad V=\Bl_C Q.
\]
Let \(E\subset V\) be the strict transform of \(H\). Let \(S=\mathbb F_1\)
be the smooth toric surface with anticanonical polygon
\[
 P=\{(x,y):-1\leq y\leq1,\ -1\leq x\leq y+1\},
\]
and let \(D\subset S\) be the toric divisor with inward primitive normal
\((0,1)\). Put \(X=V\times S\).
Define \(a,b>0\) by
\begin{align}
 &\int_{-1}^0 s(2+s)^2e^{-as}\,ds+
             \int_0^2 s(2-s)^2e^{-as}\,ds=0, \label{eq:a}\\
 &\int_{-1}^1 y(2+y)e^{-by}\,dy=0. \label{eq:b}
\end{align}
Their existence and uniqueness are proved below.

\begin{theorem}\label{thm:main}
The variety \(X\) is a smooth Fano fivefold. The quasi-monomial valuation
\(v\) with weights \((a,b)\) along the transverse divisors
\(E\times S\) and \(V\times D\) is the unique global \(H\)-minimizing
valuation on \(X\). Its value group is \(a\Z+b\Z\), of rational rank two,
so \(v\) is nondivisorial.
Its canonical first central variety is
\[
 Z\times S,\qquad
 Z=\Bl_C Q_0,\qquad
 Q_0=\{x_0^2+x_1^2+x_2^2+x_3^2=0\}\subset\mathbb P^4.
\]
This central variety is singular and hence is not isomorphic to \(X\).
It already admits the required Kähler--Ricci soliton, so the second,
polystable step changes no underlying variety.
\end{theorem}

We use the global valuative \(H\)-functional in the Han--Li/Blum--Liu--Xu--Zhuang
convention. For finite-discrepancy valuations its expression is
\[
 H(v)=A_X(v)+\log\int_{\R}e^{-\lambda}\,d\mu_v(\lambda),
\]
where \(\mu_v\) is the probability-normalized asymptotic
Duistermaat--Heckman measure of the raw valuation filtration on
\(\bigoplus_m H^0(X,-mK_X)\); see \cite[Definition~4.1]{HL}.
The scale of a valuation matters for this functional.

\section{Geometry and the established factor degeneration}

The quadric \(Q\) and its hyperplane section \(H\) are smooth.
Dependent gradients for the two equations of \(C\subset\mathbb P^3\)
would give \((i+1-\lambda)x_i=0\) for all \(i=0,\ldots,3\).
Distinct coefficients permit at most one nonzero coordinate, contradicting
\(\sum x_i^2=0\). Thus \(C\) is smooth. On \(H\simeq\mathbb P^1\times
\mathbb P^1\) it is a connected \((2,2)\)-curve of genus one:
connectedness follows from \(H^1(H,\mathcal O_H(-2,-2))=0\), and
adjunction gives its genus.

Write \(h\) for the pulled-back hyperplane class and \(F\) for the
blowup exceptional divisor. The codimension-two blowup formula gives
\(-K_V=3h-F\). The ideal \(\mathcal I_C(2)\) is globally generated by
the second quadric equation and the products of \(u\) with the five
coordinates. The graph of these generators embeds the blowup into
\(Q\times\mathbb P^N\), with second hyperplane class \(2h-F\).
Hence \(3h-F\) is ample. Also \(E\simeq H\), since \(C\) is Cartier in
\(H\).

These data put \(V\) in family \(2.23(a)\) of \cite{MW}.
Miao--Wang identify the \(\ord_E\) special degeneration, its central
variety \(Z\), and its canonical weight coordinate \(s=k/m-1\)
\cite[\S3.1, Theorems~3.1--3.2]{MW}. The weight measure has density
\[
 \rho(s)=
 \begin{cases}(2+s)^2&-1\leq s\leq0,\\
               (2-s)^2&0\leq s\leq2.
 \end{cases}
\]
Its mass is \(5\), and its positive soliton scale is exactly the \(a\)
in \eqref{eq:a}. Their weighted K-polystability theorem for
\((Z,a\xi_E)\) is an explicit input here.
The underlying ordinary family is the blowup along \(C\times\mathbb A^1\)
of
\[
 \{t^2u^2+\textstyle\sum x_i^2=0\}\subset\mathbb P^4\times\mathbb A^1 .
\]
The center is relatively smooth of codimension two; its regular-parameter
blowup charts show flatness and identify the scheme-theoretic central
fiber. The specialness and associated anticanonical grading are those
of the cited \(\ord_E\) degeneration.

The vertex \(o=[1:0:0:0:0]\) of \(Q_0\) is disjoint from \(C\).
The blowup is an isomorphism there, so \(Z\) has an ordinary
three-dimensional double point and is smooth elsewhere. Blowing up that
node resolves it with discrepancy \(1\); in particular it is klt and
Gorenstein. The same graph argument proves \(-K_Z=3h-F\) ample.

The polygon \(P\) has cyclic inward normals
\((1,0),(0,1),(-1,1),(0,-1)\), with consecutive determinants \(1\)
and all facets at distance one. Thus \(S\) is smooth Fano.
Products preserve the stated Fano and klt properties. Moreover
\(\operatorname{Sing}(Z\times S)=\{o\}\times S\), whereas \(X\) is smooth.

\section{Exact weights and irrationality}

Each left side of \eqref{eq:a}--\eqref{eq:b} is strictly decreasing:
its derivative is a negative integral of \(s^2\rho(s)e^{-ts}\)
or \(y^2(2+y)e^{-ty}\). Its value at zero is respectively
\(5/12\) or \(2/3\). A fixed negative subinterval dominates as
\(t\to+\infty\). Hence each equation has exactly one positive root.

The \(x\)-slice of \(P\) has width \(2+y\) and midpoint \(y/2\).
Consequently \eqref{eq:b} kills both components of the weighted
barycenter with weight \(e^{-by}\).
Strict positive definiteness of the covariance Hessian of
\(\log\int_Pe^{-\langle\xi,p\rangle}dp\) gives the unique toric
parameter \(\xi=(0,b)\). Wang--Zhu's theorem supplies the soliton
\cite{WZ}. Their analytic point-action parameter is the negative of
the filtration parameter used here; see \cite[Remark~2.45]{HL}.

\begin{lemma}\label{lem:irrational}
The ratio \(a/b\) is irrational.
\end{lemma}
\begin{proof}
Elementary integration, for \(t>0\), rewrites the two moment equations as
\begin{align}
 &(a^3+a^2-2a-6)e^{3a}+16ae^{2a}+4a+6=0, \label{eq:pa}\\
 &(b^2-2)e^{2b}+3b^2+4b+2=0. \label{eq:pb}
\end{align}
Here \(b^2\neq2\), since the remaining coefficient would be
\(8+4b\neq0\). If \(b\) were algebraic, \eqref{eq:pb} would make
\(e^{2b}\) algebraic, contrary to Hermite--Lindemann. Thus \(b\)
is transcendental.

Suppose \(a/b=p/q\), with positive integers \(p,q\), and write
\(r=p/q\), \(z=e^{b/q}\). Then
\[
 A(b,z)=B(b,z)=0,
\]
where
\begin{align*}
 A(T,U)&=(r^3T^3+r^2T^2-2rT-6)U^{3p}
                           +16rTU^{2p}+4rT+6,\\
 B(T,U)&=(T^2-2)U^{2q}+3T^2+4T+2.
\end{align*}
The polynomial \(R(T)=\Res_U(A,B)\in\Q[T]\) has degree at most
\(6(p+q)\), by separate homogeneity of the resultant.
Its coefficient in that degree is
\[
 \Res_U(r^3U^{3p},\,U^{2q}+3)
       =r^{6q}3^{3p}\neq0 .
\]
Indeed these are exactly the highest \(T\)-degree parts of the
two polynomials in the Sylvester determinant. Thus \(R\) is nonzero.
The common root \(z\) implies \(R(b)=0\), contradicting the
transcendence of \(b\).
\end{proof}

At the generic point of \(E\times D\), choose normal parameters \(h,w\).
The usual monomial valuation of weights \(a,b\) extends to \(\C(X)\).
Equivalently, choose toric characters \(z,w\) on \(S\) with
\(\ord_D(z)=0\), \(\ord_D(w)=1\); its Gauss formula is
\[
 v\Bigl(\sum f_{ij}z^iw^j\Bigr)
       =\min_{i,j}\{a\,\ord_E(f_{ij})+bj\}.
\]
For a fixed leading pair, distinct powers of \(z\) cannot cancel
over \(\C(E)\), so multiplication of leading terms justifies the
formula and its fraction-field extension.
All values lie in \(a\Z+b\Z\); the parameters \(h,w\) attain
\(a,b\), giving equality. Lemma~\ref{lem:irrational} makes its
rational rank two. Every \(c\,\ord_F\) has rational rank one.
The valuation has ordered rank one as a subgroup of \(\R\);
it is the \emph{rational} rank that proves nondivisoriality.

\section{The real filtration and its normalization}

Let \(A_m=H^0(V,-mK_V)\) and \(B_m=H^0(S,-mK_S)\).
Künneth gives \(R(X)_m=A_m\otimes B_m\).
For a first-factor section of \(E\)-order \(k\) and a toric
character \((i,j)\in mP\cap\Z^2\), the raw weight is
\[
                  ak+b(j+m).
\]
An \(\ord_E\)-adapted basis tensored with character lines is
adapted to \(v\). Irrationality separates distinct \((k,j)\);
different \(i\) remain linearly independent in the residue field.
Leading coefficients multiply without an additional cancellation.
Therefore, as degree-graded algebras,
\[
       \gr_v R(X)\simeq(\gr_{\ord_E}R(V))\mathbin{\#}R(S),
\]
where \(\#\) is the equal-degree Segre product.
The factor degeneration identifies \(\gr_{\ord_E}R(V)\) with the
anticanonical section ring of \(Z\), after a common Veronese if needed.
The Segre algebra is finitely generated and its Proj is \(Z\times S\)
with anticanonical polarization.

This is finite generation in the geometric real-filtration sense:
finitely many homogeneous section generators with real weights.
It is not a claim that the positive cone of the dense group
\(a\Z+b\Z\) is a finitely generated semigroup
\cite[Definition~2.2]{BLXZ}.
The ordinary \(E\) family and the toric action on \(S\), or an
adapted finite anticanonical embedding with these real diagonal
weights, realize the special real degeneration with this central ring
\cite[\S2.2, Geometric R-TC I]{HL}.
Rational directions approximate it by ordinary special degenerations;
no such rational direction is asserted to minimize \(H\).

Both transverse divisors have discrepancy one on the smooth \(X\),
so \(A_X(v)=a+b\). The canonical anticanonical weight is consequently
\[
           ak+b(j+m)-m(a+b)=a(k-m)+bj.                 \tag{*}
\]
Thus the induced central parameter is precisely
\((a\xi_E,b(0,1))\). In particular no unaccounted character twist
or sign change appears when applying the optimality criterion.

\section{Global optimality and the canonical first step}

We spell out the analytic product step to avoid assuming a tensorization
theorem for \(H\)-minimizers. The polystability of \((Z,a\xi_E)\)
from \cite{MW}, together with \cite{BLXZ,HLYTD}, gives its weak soliton.
The toric factor has the soliton specified above.
Use the Han--Li convention
\[
 e^{-\theta_\phi(\xi)}\MA_n(\phi)=e^{-\phi},\qquad
 \Ric(\omega)-\omega=-dd^c\theta_\phi(\xi),
 \quad \MA_n(\phi)=\omega^n/n!,
\]
with harmless constant normalizations of the canonical volume measure
\cite[(110)--(111), Remark~2.45]{HL}.
For factor potentials \(\phi_Z,\phi_S\), put
\(\phi=\phi_Z+\phi_S\), \(\theta=\theta_Z+\theta_S\).
The binomial expansion gives exactly
\[
 \MA_5(\phi)=\MA_3(\phi_Z)\otimes\MA_2(\phi_S).
\]
The factor equations multiply to the same equation on \(Z\times S\)
for the sum parameter. The potentials are bounded on a resolution
\cite[Corollary~31, Proposition~32]{HLYTD};
the product Monge--Ampère measure and klt canonical measure have
no mass on the analytic singular set. Hence the equality extends
from the regular locus as a weak equation. This also follows in the
finite-energy formulation of \cite{HLYTD}.

By \cite[Theorem~1.3]{BLXZ}, the central pair is weighted K-polystable,
in particular semistable.
Han--Li's criterion \cite[Theorem~5.3]{HL}, with their comparison
of infima \cite[Theorem~1.5]{HL}, now proves that \(v\) minimizes
\(H\) among \emph{all} valuations on the original \(X\).
The special real degeneration, its klt anticanonical central ring,
finite discrepancy, and exact induced parameter were checked above;
no restriction to product-form competitors is made.
Blum--Liu--Xu--Zhuang's absolute uniqueness theorem
\cite[Theorem~1.2]{BLXZ} identifies this valuation and
\(\Proj\gr_vR(X)\) with the canonical first step.
The central pair is already polystable, so its second step changes
no underlying variety. The smooth/singular distinction and
Lemma~\ref{lem:irrational} complete Theorem~\ref{thm:main}.

\section{Scope and verification}

This is an explicit construction from the cited general theorems and
Miao--Wang's established threefold example. We do not claim historical
firstness, a new general optimal-degeneration theorem, a
three-dimensional nondivisorial example, or an exhaustive priority
determination. The statement resolves the specified construction
request in dimension five without a product exclusion.
The accompanying archive preserves the original proof and corrections,
primary-source fingerprints, independent adversarial analyses, and
exact auxiliary checks. Numerical brackets and finitely many
resultants do not prove irrationality; the universal argument above does.
Successful software checks do not formally verify the geometric or
analytic theorem inputs.

AI tools were used extensively in developing, checking, and preparing this
work. The reported adversarial reviews were conducted by AI agents.
This preprint is unrefereed; independent human peer review and formal
verification have not been completed.

\begin{thebibliography}{9}
\bibitem{OWR}
H.~Blum, Degenerations of unstable Fano varieties and Kahler--Ricci soliton,
in \emph{Toric Geometry}, \emph{Oberwolfach Rep.} 19 (2022), no.~2,
Report 17/2022, pp.~888--889; published online 14 April 2023.
\href{https://doi.org/10.4171/OWR/2022/17}{doi:10.4171/OWR/2022/17}.
\bibitem{MW}
M.~Miao and L.~Wang, \emph{Optimal Degenerations of K-unstable Fano
threefolds}, first posted 25 January 2024; cited edition
arXiv:2401.13999v2, 13 April 2026.
\href{https://arxiv.org/abs/2401.13999v2}{arXiv:2401.13999v2}.
\bibitem{LL}
Y.~Li and Z.~Li, Equivariant \(\mathbb R\)-test configurations and
semistable limits of \(\mathbb Q\)-Fano group compactifications,
\emph{Peking Math. J.} 6 (2023), 559--607.
\href{https://doi.org/10.1007/s42543-022-00054-0}{doi:10.1007/s42543-022-00054-0}.
The inspected example text is
\href{https://arxiv.org/abs/2103.06439v6}{arXiv:2103.06439v6}.
\bibitem{WZ}
X.-J.~Wang and X.~Zhu, Kähler--Ricci solitons on toric manifolds with
positive first Chern class, \emph{Adv. Math.} 188 (2004), 87--103.
\href{https://doi.org/10.1016/j.aim.2003.09.009}{doi:10.1016/j.aim.2003.09.009}.
\bibitem{HL}
J.~Han and C.~Li, Algebraic uniqueness of Kähler--Ricci flow limits
and optimal degenerations of Fano varieties,
\emph{Geom. Topol.} 28 (2024), 539--592.
The cited theorem numbering is that of
\href{https://arxiv.org/abs/2009.01010v1}{arXiv:2009.01010v1}.
\bibitem{HLYTD}
J.~Han and C.~Li, On the Yau--Tian--Donaldson conjecture for
generalized Kähler--Ricci soliton equations,
\emph{Comm. Pure Appl. Math.} 76 (2023), 1793--1867.
\href{https://doi.org/10.1002/cpa.22053}{doi:10.1002/cpa.22053}.
The cited corollary/proposition locators refer to the
\href{https://sites.math.rutgers.edu/~cl1412/papers/g-soliton.pdf}{inspected author manuscript}.
\bibitem{BLXZ}
H.~Blum, Y.~Liu, C.~Xu and Z.~Zhuang, The existence of the
Kähler--Ricci soliton degeneration,
\emph{Forum Math. Pi} 11 (2023), e9.
\href{https://doi.org/10.1017/fmp.2023.5}{doi:10.1017/fmp.2023.5}.
The inspected preprint is
\href{https://arxiv.org/abs/2103.15278v2}{arXiv:2103.15278v2}.
\end{thebibliography}
\end{document}
